% Generated from locale/tr/content/intuitionistic-logic/introduction/syntax.tex; do not hand-edit.


\olfileid[tr]{int}{int}{syn}

\olsection{Sezgici Mantığın Sözdizimi}

Sezgici mantığın sözdizimi, önerme mantığınınkiyle aynıdır. Klasik önerme
mantığında bağlaçları başka bağlaçlar aracılığıyla tanımlamak mümkündür;
örneğin $!A \lif !B$, $\lnot !A \lor !B$ olarak ya da $!A \lor !B$,
$\lnot(\lnot !A \land \lnot !B)$ olarak tanımlanabilir. Bu nedenle klasik
mantık sunumları, bazı bağlaçları çoğu kez bu tanımların kısaltmaları olarak
tanıtır. Sezgici mantıkta iki istisna dışında durum böyle değildir: $\lnot !A$,
$!A \lif \lfalse$ ifadesinin kısaltması olarak tanımlanabilir ve çoğu kez de
öyle tanımlanır. Elbette bu durumda $\lfalse$'ın kendisi tanımlanmamalıdır!{}
Ayrıca $!A \liff !B$, klasik mantıkta olduğu gibi $(!A \lif !B) \land (!B
\lif !A)$ olarak tanımlanabilir.

Önermesel sezgici mantığın !!^{formula}s{}i, \emph{!!{propositional variable}s}
ile önerme sabiti~$\lfalse$'tan \emph{mantıksal bağlaçlar} kullanılarak
kurulur. Elimizde şunlar vardır:

\begin{enumerate}
\item !!^a{denumerable} !!{propositional variable} kümesi~$\PVar$: $\Obj p_0$,
  $\Obj p_1$, \dots
\item !!{falsity} için önerme sabiti~$\lfalse$.
\item Mantıksal bağlaçlar: $\land$ (ve bağlacı), $\lor$ (veya bağlacı),
  $\lif$ (koşul).
\item Noktalama işaretleri: (, ) ve virgül.
\end{enumerate}

\begin{defn}[Formül]
\ollabel{defn:formulas} $\Frm[L_0]$ kümesi, önermesel sezgici mantığın
\emph{!!{formula}s{}den} oluşur ve tümevarımlı olarak şöyle tanımlanır:
\begin{enumerate}
\item $\lfalse$ atomik !!{formula}{}dür.

\item Her !!{propositional variable}~$\Obj p_i$ atomik !!{formula}{}dür.

\item $!A$ ve $!B$ birer !!{formula} ise $(!A \land !B)$ de
  !!a{formula}{}dür.

\item $!A$ ve $!B$ birer !!{formula} ise $(!A \lor !B)$ de
  !!a{formula}{}dür.

\item $!A$ ve $!B$ birer !!{formula} ise $(!A \lif !B)$ de
  !!a{formula}{}dür.

\tagitem{limitClause}{Başka hiçbir şey !!a{formula} değildir.}{}
\end{enumerate}
\end{defn}

Yukarıda tanıtılan ilkel bağlaçlara ek olarak şu \emph{tanımlı} simgeleri de
kullanırız: $\lnot$ (değilleme) ve $\liff$ (!!{biconditional}). Tanımlı
işleçler kullanılarak kurulan formüller şöyle anlaşılmalıdır:
\begin{enumerate}
\item $\lnot !A$, $!A \lif \lfalse$ ifadesinin kısaltmasıdır.

\item $!A \liff !B$, $(!A \lif !B) \land (!B \lif !A)$ ifadesinin
  kısaltmasıdır.
\end{enumerate}

$\lnot$ resmen bir kısaltma olarak ele alınsa da tanımlarda kimi zaman~$\lnot$
için ilkelmiş gibi açık kurallar ve koşullar vereceğiz. Bunu daha çok alıştırma
soruları yazabilmek için yapacağız.

